\documentclass{article}
\usepackage[T1]{fontenc}
\usepackage[utf8]{inputenc}

\begin{document}

Ovo je uvodni deo teksta sa jednim osnovnim citatom \cite{knjigaA}.

U sledecoj recenici isti citat se ponavlja \cite{knjigaA}, pa ovde ne bi trebalo da nastane nova fusnota.

Sada navodimo grupni citat \cite{clanakB,radC,konfD}, pa ocekujemo jednu zajednicku fusnotu sa tri bibliografske jedinice razdvojene tacka-zarezom.

Ovde imamo potpuno nepoznat kljuc \cite{nepoznat1}.

U ovoj recenici imamo mesoviti slucaj \cite{knjigaA,nepoznat2,knjigaE}, gde je jedan citat vec ranije koriscen, jedan je nepoznat, a jedan se pojavljuje prvi put.
Onda opet \cite{knjigaA,nepoznat2,knjigaE}.
Imamo i podskup \cite{knjigaA,nepoznat2}.
Kao i nadskup \cite{knjigaA,nepoznat2,knjigaE,radF}.

Sada proveravamo citat koji ima prazna polja \cite{nepotpunZapis}.

Ovde koristimo dva nova citata zajedno \cite{radF,knjigaG}, pa bi trebalo da budu u istoj fusnoti.

Ovde imamo jos jedan grupni primer \cite{knjigaA,clanakB,radF}, gde su svi citati vec ranije pomenuti, ali se sada pojavljuju zajedno, pa bi opet trebalo da budu ispisani u jednoj zajednickoj fusnoti.

Sada navodimo citat sa vise polja \cite{detaljanRad}.

Na kraju ponavljamo jedan vec koriscen pojedinacni citat \cite{detaljanRad}.

\bibliographystyle{plain}
\bibliography{literatura9}


\end{document}